%% 10-limitations
\section{Limitations and Outlook}
\label{sec:limitations}

\paragraph{No quantum speedup in kernel evaluation.}
Restating the point of \Cref{sec:introduction} where a reader will look for it:
$K(x,y) = f(\lcadepth(x,y))$ is computable classically in time linear in the tree
depth, so nothing here accelerates kernel evaluation. The obstructions of
\Cref{thm:product-nogo,thm:dimension-bound} constrain quantum encoding design, and the
construction of \Cref{thm:path-states} is useful as an input stage for deeper quantum
models, where the closed form does not survive subsequent layers. Neither is a speedup
claim.

\paragraph{Multiple inheritance is discarded.}
All three real hierarchies are directed acyclic graphs, and the longest-path reduction
of \Cref{subsec:setup} throws away every hypernym but one. A synset with two genuinely
distinct parents is misrepresented. Extending the construction to DAGs is not a matter
of a better reduction: the ancestor sets of two leaves in a DAG need not be nested, and
\cref{eq:overlap} --- the entire content of \Cref{thm:path-states} --- depends on that
nesting.

\paragraph{Truncation.}
Comparing against $2^{\text{height}}$-dimensional baselines forced a depth-$8$
truncation. The path-state encoding itself has no such limit --- \Cref{tab:resources}
gives $17$ qubits for the untruncated WordNet hierarchy --- so the truncation is a
constraint of the comparison, not of the method. We have not measured the baselines at
full depth, because they cannot be represented there.

\paragraph{The dimension bound is loose for flat profiles.}
\Cref{thm:dimension-bound} is strong when the profile decays quickly and weak when it
does not, and \Cref{rem:s-greater-one} shows the hypothesis $s > 1$ is necessary rather
than convenient. For slowly decaying profiles the bound does not rule out narrow
encodings, and we do not know whether narrow encodings actually exist in that regime.

\paragraph{Simulation only.}
Everything here is statevector simulation. \Cref{prop:depolarising} covers the global
depolarising channel analytically, and \Cref{subsec:e5} gives a shot budget of
$\ShotBudget$ per kernel entry, but coherent errors, readout error and device-specific
noise are untested. The shot budget is also per entry, so a full $m \times m$ kernel
matrix costs $O(m^{2})$ times that, which is the practical barrier to using this at
scale on hardware.

\paragraph{E2 is parity evidence, not a win.}
As reported in \Cref{subsec:e2}, the path state loses on classification accuracy to
both a random product map and to v-PuNNs. The paper's claims rest on
\Cref{thm:product-nogo,thm:dimension-bound,thm:path-states} and on E1, not on E2.

\paragraph{Outlook.}
Three directions seem worth pursuing. The most immediate is using the path state as the
input layer of a deeper variational model, where \Cref{subsec:e4} suggests the
interesting regime is small distortion rather than large. The second is the $p$-adic
connection: the encoding lives in the state space of quantum walks on the $p$-adic tree
\citep{ZunigaGalindo2024QuantumWalks}, and relating the profile $f$ to the spectrum of
a walk generator would give a dynamical reading of what is currently a static
construction. The third is gate synthesis by $p$-adic descent, where the ultrametric
structure would organise the search rather than the data; that is a longer project and
we mention it only to place the present work.
